\exposetitle{IV}{Topologies and sheaves}
\label{IV}

\begin{center}
\textit{by M.~Demazure}%
{\renewcommand{\thefootnote}{(*)}\footnote{This text develops the substance
of two oral expos\'es by A.~Grothendieck, completing the latter on several
important points, which had been passed over in silence or barely touched
upon.}}
\end{center}

{\renewcommand{\thefootnote}{(0)}\nde{: Version of 13/10/2024}}%
This expos\'e is intended to acquaint the reader with the essentials of the
language of topologies and of sheaves (without cohomology), particularly
convenient in questions of passage to the quotient (among others).

The first three paragraphs develop the language of passage to the quotient.
The fourth, which is the central part, is the exposition of the theory of
sheaves, oriented principally towards the application to questions of
quotients; the fifth is an application to passage to the quotient in groups
and to principal homogeneous bundles. The last paragraph concerns more
specifically the category of schemes, and defines various useful topologies
on this category.

The reader will find it useful to refer to \textup{[AS]}, \textup{[MA]},
\textup{[D]}, and SGA~4; to \textup{[D]} especially concerning the
applications of topologies to descent theory, and to SGA~4 for questions of
universes (particularly ill-treated in this expos\'e).

\sgasection{1}{Universal effective epimorphisms}
\label{IV.sec.1}

In the sequel of this expos\'e, one assumes fixed a category~\(\mathcal{C}\).

\begin{definition}{1.1}
\label{IV.1.1}
A morphism \(u\colon T\to S\) is called an \emph{epimorphism} if, for every
object~\(X\), the corresponding map
\[
X(S)=\Hom(S,X)\longrightarrow X(T)=\Hom(T,X)
\]
is injective%
{\renewcommand{\thefootnote}{(1)}\nde{: that is, if \(u\) is
right-cancellable.}}.
One says that \(u\) is a \emph{universal epimorphism} if for every morphism
\(S'\to S\), the fiber product \(T'=T\times_{S}S'\) exists, and
\(u'\colon T'\to S'\) is an epimorphism.
\end{definition}

\begin{definition}{1.2}
\label{IV.1.2}
A diagram
\[
\xymatrix@C=3em{
A \ar[r]^{u}
  & B \ar@<0.5ex>[r]^{v_{1}} \ar@<-0.5ex>[r]_{v_{2}}
  & C
}
\]
of maps of sets is said to be \emph{exact} if \(u\) is injective and if its
image consists of the elements \(b\) of~\(B\) such that
\(v_{1}(b)=v_{2}(b)\). A diagram of the same type in~\(\mathcal{C}\) is said
to be \emph{exact} if for every object \(X\) of~\(\mathcal{C}\), the
corresponding diagram of sets
\[
\xymatrix@C=3em{
A(X) \ar[r]
  & B(X) \ar@<0.5ex>[r] \ar@<-0.5ex>[r]
  & C(X)
}
\]
is exact; one then also says that \(u\) makes \(A\) a \emph{kernel} of the
pair of arrows \((v_{1},v_{2})\). Dually, a diagram
\[
\xymatrix@C=3em{
C \ar@<0.5ex>[r]^{v_{1}} \ar@<-0.5ex>[r]_{v_{2}}
  & B \ar[r]^{u}
  & A
}
\]
in~\(\mathcal{C}\) is said to be \emph{exact} if it is exact as a diagram in
the opposite category \(\mathcal{C}^{\circ}\), \ie if for every object \(X\)
of~\(\mathcal{C}\), the corresponding diagram of sets
\[
\xymatrix@C=3em{
X(A) \ar[r]
  & X(B) \ar@<0.5ex>[r] \ar@<-0.5ex>[r]
  & X(C)
}
\]
is exact.%
{\renewcommand{\thefootnote}{(2)}\nde{: This implies, in particular, that
\(u\) is an epimorphism.}}
One also says that \(u\) makes \(A\) a \emph{cokernel} of the pair of arrows
\((v_{1},v_{2})\).
\end{definition}

\begin{definition}{1.3}
\label{IV.1.3}
A morphism \(u\colon T\to S\) is called an \emph{effective epimorphism} if
the fiber square \(T\times_{S}T\) exists, and if the diagram
\[
\xymatrix@C=3.2em{
T\times_{S}T
  \ar@<0.5ex>[r]^{\mathrm{pr}_{1}}
  \ar@<-0.5ex>[r]_{\mathrm{pr}_{2}}
  & T \ar[r]^{u}
  & S
}
\]
is exact, \ie if \(u\) makes \(S\) a cokernel of
\((\mathrm{pr}_{1},\mathrm{pr}_{2})\). One says that \(u\) is a
\emph{universal effective epimorphism} if for every morphism \(S'\to S\),
the fiber product \(T'=T\times_{S}S'\) exists, and the morphism
\(u'\colon T'\to S'\) is an effective epimorphism.
\end{definition}

One evidently has the implications:
\[
\xymatrix@C=3.2em@R=2.8em{
\mbox{\shortstack{universal effective\\epimorphism}} \ar[r] \ar[d]
  & \mbox{\shortstack{effective\\epimorphism}} \ar[d] \\
\mbox{\shortstack{universal\\epimorphism}} \ar[r]
  & \mbox{epimorphism}\,,
}
\]
but in general no other implication is valid.%
{\renewcommand{\thefootnote}{(3)}\nde{: For example, if
\(\mathcal{C}=\Sch\) is the category of schemes, one easily sees that every
universal epimorphism is surjective. Let
\(T=\coprod_{p\ \mathrm{prime}}\Spec(\FF_{p})\) and \(S=\Spec(\ZZ)\); then
the morphism \(u\colon T\to S\) is an epimorphism which is \emph{not
universal}. On the other hand, one sees that \(T\times_{S}T\) is identified
with~\(T\), so that the two projections \(T\times_{S}T\rightrightarrows T\)
coincide; as \(\mathrm{id}_{T}\) does not descend to a morphism \(S\to T\),
this shows that \(u\) is \emph{not} an effective epimorphism.}}

\begin{definition}{1.4.0}
\label{IV.1.4.0}
{\renewcommand{\thefootnote}{(4)}\nde{: One has added the numbering~1.4.0,
for later references.}}%
One ``recalls'' that a morphism \(u\colon T\to S\) is said to be
\emph{quarrable} if for every morphism \(S'\to S\), the fiber product
\(T\times_{S}S'\) exists.
\end{definition}

\begin{lemma}{1.4}
\label{IV.1.4}
Consider morphisms \(U\xrightarrow{v}T\xrightarrow{u}S\). Then

a)~ \(u\), \(v\) epimorphisms \(\Rightarrow\) \(uv\) epimorphism \(\Rightarrow\)
\(u\) epimorphism,

b)~ \(u\), \(v\) universal epimorphisms \(\Rightarrow\) \(uv\) universal
epimorphism and \(u\) quarrable \(\Rightarrow\) \(u\) universal epimorphism.
\end{lemma}

Lemma~1.4 is trivial from the definitions. One concludes from it:

\begin{corollary}{1.5}
\label{IV.1.5}
Let \(u\colon X\to Y\) and \(u'\colon X'\to Y'\) be universal epimorphisms,
such that \(Y\times Y'\) exists, then \(X\times X'\) exists and
\(u\times u'\colon X\times X'\to Y\times Y'\) is a universal epimorphism.
\end{corollary}

Let us also note:

\begin{definition}{1.6.0}
\label{IV.1.6.0}
{\renewcommand{\thefootnote}{(5)}\nde{: One has added the numbering~1.6.0,
for later references.}}%
One says that an object \(S\) of~\(\mathcal{C}\) is \emph{quarrable} if its
product by every object of~\(\mathcal{C}\) exists. (If \(\mathcal{C}\) has a
final object~\(e\), this amounts to saying that the morphism \(S\to e\) is
quarrable, \cf~1.4.0.)
\end{definition}

\begin{lemma}{1.6}
\label{IV.1.6}
Let \(u\colon X\to Y\) be a morphism in \(\mathcal{C}_{/S}\); for it to be
an epimorphism (\resp a universal epimorphism, \resp an effective
epimorphism, \resp a universal effective epimorphism), it suffices that the
corresponding morphism in~\(\mathcal{C}\) be one, and this is also necessary
if one assumes that \(S\) is a quarrable object of~\(\mathcal{C}\).
\end{lemma}

\begin{proof}
Immediate, left to the reader. One uses the hypothesis ``\(S\) quarrable''
in order to interpret the \(\mathcal{C}\)-morphisms from an object \(Y\)
of \(\mathcal{C}_{/S}\) into an object \(Z\) of~\(\mathcal{C}\), as being the
\(\mathcal{C}_{/S}\)-morphisms from \(Y\) into \(Z\times S\).
\end{proof}

\begin{lemma}{1.7}
\label{IV.1.7}
With the notations of~1.4: \(u\), \(v\) effective epimorphisms and \(v\)
universal epimorphism \(\Rightarrow\) \(uv\) effective epimorphism.
\end{lemma}

To see this, one considers the diagram
\[
\xymatrix@C=3.6em@R=2.6em{
S
  & T \ar[l]_{u}
  & T\times_{S}T \ar@<0.45ex>[l] \ar@<-0.45ex>[l]
\\
  & U \ar[u]_{v}
  & U\times_{S}U \ar[u]_{v\times_{S}v}
      \ar@<0.45ex>[l] \ar@<-0.45ex>[l]
\\
  & U\times_{T}U \ar@<0.45ex>[u] \ar@<-0.45ex>[u] \ar[ur]
  &
}
\,.
\]
One notes that by hypothesis, row~1 and column~1 are exact, and that by
virtue of~1.5 and~1.6, \(v\times_{S}v\) is an epimorphism (\(v\) being a
\emph{universal} epimorphism). The conclusion follows from this by an
evident diagram-chasing: if an element of \(X(U)\) has the same images in
\(X(U\times_{S}U)\), it \emph{a fortiori} has the same images in
\(X(U\times_{T}U)\), hence comes from an element of \(X(T)\) since column~1
is exact. As row~1 is exact, it suffices to check that the element under
consideration has the same images in \(X(T\times_{S}T)\), and as
\(v\times_{S}v\) is an epimorphism, it suffices to check that the images in
\(X(U\times_{S}U)\) are the same, which is indeed the case.

\begin{proposition}{1.8}
\label{IV.1.8}
Consider morphisms \(U\xrightarrow{v}T\xrightarrow{u}S\). Then \(u\), \(v\)
universal effective epimorphisms \(\Rightarrow\) \(uv\) universal effective
epimorphism and \(u\) quarrable \(\Rightarrow\) \(u\) universal effective
epimorphism.
\end{proposition}

The first implication follows at once from~1.7. For the second, one looks
at the diagram (of ``bisimplicial'' type):
{\footnotesize\[
\xymatrix@C=0.55em@R=2.2em{
S
  & T \ar[l]_{u}
  & T\times_{S}T \ar@<0.4ex>[l] \ar@<-0.4ex>[l]
\\
U \ar[u]_{uv}
  & U\times_{S}T \ar[l] \ar[u]
  & U\times_{S}T\times_{S}T
      \ar@<0.4ex>[l] \ar@<-0.4ex>[l] \ar[u]
\\
U\times_{S}U \ar@<0.4ex>[u] \ar@<-0.4ex>[u]
  & U\times_{S}U\times_{S}T \ar[l]
      \ar@<0.4ex>[u] \ar@<-0.4ex>[u]
  & U\times_{S}U\times_{S}T\times_{S}T
      \ar@<0.4ex>[l] \ar@<-0.4ex>[l]
      \ar@<0.4ex>[u] \ar@<-0.4ex>[u]
}
\,.
\]}
Columns~1, 2, 3 are exact by virtue of the hypothesis ``\(uv\) universal
effective epimorphism'', row~2 is exact, for \(U\times_{S}T\to U\) is an
effective epimorphism (for it has a section over~\(U\)), and the same holds
for row~3 (same reason). An evident diagram-chasing then shows that row~1
is exact, \ie \(u\) is an effective epimorphism. As the hypotheses made are
invariant under an arbitrary base change \(S'\to S\), it follows that \(u\)
is even a universal effective epimorphism.

\begin{corollary}{1.9}
\label{IV.1.9}
Let \(u\colon X\to Y\) and \(u'\colon X'\to Y'\) be universal effective
epimorphisms, such that \(Y\times Y'\) exists; then \(X\times X'\) exists
and \(u\times u'\colon X\times X'\to Y\times Y'\) is a universal effective
epimorphism.
\end{corollary}

\begin{proof}
As for~1.5, by the diagram
\[
\xymatrix@C=3.2em@R=2.4em{
  & Y'
  & X' \ar[l]_{u'}
\\
X \ar[d]_{u}
  & X\times Y' \ar[l] \ar[u] \ar[d]
  & X\times X' \ar[l] \ar[u] \ar[dl]^{u\times u'}
\\
Y
  & Y\times Y' \ar[l]
  &
}
\,.
\]
\end{proof}

\begin{corollary}{1.10}
\label{IV.1.10}
Consider a quarrable morphism \(u\colon T\to S\), and a base change morphism
\(S'\to S\), which is a universal effective epimorphism. For \(u\) to be a
universal effective epimorphism, it is necessary and sufficient that
\(u'\colon T'=T\times_{S}S'\to S'\) be so:
\[
\xymatrix@C=3em@R=2.4em{
T \ar[d]_{u}
  & T' \ar[l]_{v'} \ar[d]^{u'}
\\
S
  & S' \ar[l]_{v}
}
\,.
\]
\end{corollary}

Only the ``it suffices'' requires a proof. Now if \(u'\) is a universal
effective epimorphism, the same holds for \(vu'\) thanks to~1.8, and since
\(vu'=uv'\), one concludes by~1.8 that \(u\) is a universal effective
epimorphism.

\begin{remark}{1.11}
\label{IV.1.11}
The same reasoning shows that in~1.10 one may replace ``universal effective
epimorphism'' by ``universal epimorphism'' or ``universal and effective
epimorphism'', or simply by ``epimorphism'' (and in this last case, the
hypothesis ``\(u\) quarrable'' is evidently useless).
\end{remark}

In the proof of~1.8 we have used the following result, which deserves to be
spelled out:

\begin{proposition}{1.12}
\label{IV.1.12}
Let \(u\colon T\to S\) be a morphism which admits a section. Then \(u\) is
an epimorphism, and if \(T\times_{S}T\) exists, it is an effective
epimorphism, and a universal effective epimorphism if moreover \(u\) is
quarrable.
\end{proposition}

The first assertion is contained in~1.4~a), and the third will follow at
once from the second, which it will therefore suffice to establish. In fact
one has a stronger conclusion: for every functor
\(F\colon\mathcal{C}^{\circ}\to\Ens\) (not necessarily representable), the
diagram of sets
\[
\xymatrix@C=3em{
F(S) \ar[r]
  & F(T) \ar@<0.5ex>[r] \ar@<-0.5ex>[r]
  & F(T\times_{S}T)
}
\]
is exact. This may be considered as a particular case of the formalism of
\v{C}ech cohomology (in dimension~0!) which we content ourselves with
recalling here. Let us simply assume that \(T\times_{S}T\) exists; one then
sets
\[
\mathrm{H}^{0}(T/S,F)
=\Ker\bigl(F(T)\rightrightarrows F(T\times_{S}T)\bigr).
\]
One may regard \(\mathrm{H}^{0}(T/S,F)\) in an evident way as a
contravariant functor in the argument \(T\) varying in \(\mathcal{C}_{/S}\),
every \(S\)-morphism \(T'\to T\) defining a map
\begin{equation}
\mathrm{H}^{0}(T/S,F)\longrightarrow\mathrm{H}^{0}(T'/S,F).
\tag{+}
\end{equation}
Let us fix \(T\) and \(T'\) in \(\mathcal{C}_{/S}\). A well-known computation
shows that if there exists an \(S\)-morphism from \(T'\) into~\(T\), the
corresponding map \((+)\) is in fact independent of the choice of this
morphism%
{\renewcommand{\thefootnote}{(6)}\nde{: This is the following argument,
communicated by M.~Demazure. Let \(f,g\colon T'\to T\), and let
\(\varphi\colon T'\to T\times_{S}T\) be the morphism with components \(f\)
and~\(g\), whence \(p_{1}\circ\varphi=f\) and \(p_{2}\circ\varphi=g\). Then
\(F(\varphi)\colon F(T\times_{S}T)\to F(T')\) satisfies
\(F(\varphi)\circ F(p_{1})=F(f)\) and \(F(\varphi)\circ F(p_{2})=F(g)\). Now,
for every \(x\in\mathrm{H}^{0}(T/S,F)\), one has
\(F(p_{1})(x)=F(p_{2})(x)\). Therefore, applying \(F(\varphi)\) to both
sides, one obtains \(F(f)(x)=F(g)(x)\), which shows that \(f\) and \(g\)
induce the same morphism.}},
so that \(\mathrm{H}^{0}(T/S,F)\) may be regarded as a functor on the
category associated with the set \(\Ob\mathcal{C}_{/S}\) preordered by the
relation of ``domination'' (\(T'\) dominates \(T\) if there exists an
\(S\)-morphism from \(T'\) into~\(T\)). In particular, if \(T\) and \(T'\) are
isomorphic in this latter category, \ie if each dominates the other, then
\((+)\) is an isomorphism of sets. This applies in particular to the case
where \(T'\) is the final object of \(\mathcal{C}_{/S}\), \ie essentially
\(S\) itself; in any case \(T\) dominates \(T'=S\), and the converse is true
precisely if \(T/S\) has a section. This establishes~1.12 in the strengthened
form announced.

\begin{remark}{1.13}
\label{IV.1.13}
For various applications, the notions introduced in the present expos\'e,
and the results stated, are to be developed more generally relative to a
\emph{family} of morphisms \(u_{i}\colon T_{i}\to S\) with the same target
(instead of a single morphism \(u\colon T\to S\)). Thus, such a family will
be said to be \emph{epimorphic} if for every object \(X\) of~\(\mathcal{C}\),
the corresponding map
\[
X(S)\longrightarrow\prod_{i}X(T_{i})
\]
is injective, and one likewise introduces the notion of effective
epimorphic family and the ``universal'' variants of these notions. We shall
admit if need be, in the sequel, that the results of the present expos\'e
extend to this more general situation.
\end{remark}

\begin{remark}{1.14}
\label{IV.1.14}
Every morphism which is at the same time a monomorphism and an effective
epimorphism is an isomorphism. Indeed, in the notations of~1.3, the fact
that \(T\to S\) is a monomorphism implies that the two morphisms
\[
\mathrm{pr}_{1},\mathrm{pr}_{2}\colon
T\times_{S}T\rightrightarrows T
\]
are equal (and are isomorphisms). Now a diagram
\[
\xymatrix@C=3em{
C \ar@<0.5ex>[r]^{v} \ar@<-0.5ex>[r]_{v}
  & B \ar[r]^{u}
  & A
}
\]
is exact only if \(u\) is an isomorphism, as follows immediately from the
definition.%
{\renewcommand{\thefootnote}{(7)}\nde{: One will note that a monomorphism
which is an epimorphism is not necessarily an isomorphism. For example, in
\(\mathcal{C}=\Sch\), the morphism
\(\Spec(\FF_{p})\coprod\Spec(\ZZ[1/p])\to\Spec(\ZZ)\) is a monomorphism and
a surjective epimorphism, but is not an isomorphism.}}
\end{remark}

\sgasection{2}{Descent morphisms}
\label{IV.sec.2}

Let us recall the following definitions:

\begin{definition}{2.1}
\label{IV.2.1}
Let \(f\colon S'\to S\) be a morphism such that \(S''=S'\times_{S}S'\)
exists, and let \(u'\colon X'\to S'\) be an object over~\(S'\). One calls
\emph{gluing datum} on \(X'/S'\), relative to~\(f\), an
\(S''\)-isomorphism
\[
c\colon X''_{1}\isomto X''_{2}
\]
where \(X''_{i}\) (\(i=1,2\)) denotes the inverse image (assumed to exist)
of \(X'/S'\) by the projection \(\mathrm{pr}_{i}\colon S''\to S'\). One says
that the gluing datum \(c\) is a \emph{descent datum} if it satisfies the
``cocycle condition''
\[
\mathrm{pr}_{3,1}^{*}(c)
=\mathrm{pr}_{3,2}^{*}(c)\,\mathrm{pr}_{2,1}^{*}(c)
\]
where \(\mathrm{pr}_{i,j}\) (\(1\leqslant j<i\leqslant 3\)) are the
canonical projections of \(S'''=S'\times_{S}S'\times_{S}S'\) into~\(S''\)
(N.B.\ one now assumes that \(S'''\) also exists), where
\(\mathrm{pr}_{i,j}^{*}(c)\) is the inverse image of~\(c\), considered as an
\(S'''\)-morphism from \(X'''_{j}\) into \(X'''_{i}\), and where for every
integer \(k\) between \(1\) and~\(3\), \(X'''_{k}\) denotes the inverse image
(assumed to exist) of \(X'/S'\) by the projection of index~\(k\),
\(q_{k}\colon S'''\to S'\).
\end{definition}
